\documentclass{article}
\usepackage[utf8]{inputenc}
\usepackage[ngerman]{babel}
\usepackage{amsmath, amssymb}

\begin{document}

\section*{Aufgabe 2}

\subsection*{a)}
Behauptung: 

Für beliebige Mengen $A, B, C$ gilt 
 $(A \cup B) \cap C = A \cup (B \cap C)$

\subsubsection*{Gewählte Mengen:}
\begin{align*}
A &= \{1,2\} \\
B &= \{2,3\} \\
C &= \{1,3\}
\end{align*}

\subsubsection*{Linke Seite: $(A \cup B) \cap C$}
\begin{align*}
	(A \cup B) \cap C &= (\{1,2\} \cup \{2,3\}) \cap \{1,3\} \\
	&= \{1,3\}
\end{align*}

\subsubsection*{Rechte Seite: $A \cup (B \cap C)$}
\begin{align*}
	A \cup (B \cap C) &= \{1,2\} \cup (\{2,3\} \cap \{1,3\}) \\
	&= \{1,2,3\}
\end{align*}

\subsubsection*{Vergleich:}
$$\{1, 3\} \neq \{1, 2, 3\}$$

q.e.d.

\subsection*{b)}
Behauptung: 

Für beliebige Mengen $A, B$ gilt 
$\mathcal{P}(A \cup B) = \mathcal{P}(A) \cup \mathcal{P}(B)$

\subsubsection*{Gewählte Mengen:}
\begin{align*}
A &= \{1,2\} \\
B &= \{2,3\}
\end{align*}

\subsubsection*{Linke Seite: $\mathcal{P}(A \cup B)$}
\begin{align*}
	\mathcal{P}(A \cup B) &= \mathcal{P}(\{1,2\} \cup \{2,3\}) \\
												&= \{\emptyset, \{1\}, \{2\}, \{3\}, \{1,2\}, \{1,3\}, \{2,3\}, \{1,2,3\}\}
\end{align*}

\subsubsection*{Rechte Seite: $\mathcal{P}(A) \cup \mathcal{P}(B)$}
\begin{align*}
	\mathcal{P}(A) \cup \mathcal{P}(B) &= \mathcal{P}(\{1,2\}) \cup \mathcal{P}(\{2,3\}) \\
																		 &= \{\emptyset, \{1\}, \{2\}, \{3\}, \{1,2\}, \{2,3\}\}
\end{align*}

\subsubsection*{Vergleich:}
$$\{\emptyset, \{1\}, \{2\}, \{3\}, \{1,2\}, \{1,3\}, \{2,3\}, \{1,2,3\}\}
 \neq \{\emptyset, \{1\}, \{2\}, \{3\}, \{1,2\}, \{2,3\}\}$$

q.e.d.

\subsection*{c)}

\subsection*{d)}
Behauptung: 

Für beliebige Mengen $A, B, C, D$ gilt 
$(A \setminus B) \times (C \setminus D) = (A \times C) \setminus (B \times D)$.

\subsubsection*{Gewählte Mengen:}
\begin{align*}
A &= \{1, 2\} \\
B &= \{2\} \\
C &= \{3, 4\} \\
D &= \{4\}
\end{align*}

\subsubsection*{Linke Seite: $(A \setminus B) \times (C \setminus D)$}
\begin{align*}
A \setminus B &= \{1, 2\} \setminus \{2\} = \{1\} \\
C \setminus D &= \{3, 4\} \setminus \{4\} = \{3\} \\
(A \setminus B) \times (C \setminus D) &= \{1\} \times \{3\} = \{(1, 3)\}
\end{align*}

\subsubsection*{Rechte Seite: $(A \times C) \setminus (B \times D)$}
\begin{align*}
A \times C &= \{1, 2\} \times \{3, 4\} = \{(1, 3), (1, 4), (2, 3), (2, 4)\} \\
B \times D &= \{2\} \times \{4\} = \{(2, 4)\} \\
(A \times C) \setminus (B \times D) &= \{(1, 3), (1, 4), (2, 3), (2, 4)\} \setminus \{(2, 4)\} \\
&= \{(1, 3), (1, 4), (2, 3)\}
\end{align*}

\subsubsection*{Vergleich:}
$$\{(1, 3)\} \neq \{(1, 3), (1, 4), (2, 3)\}$$

q.e.d.

\section*{Aufgabe 3}
Der Term besagt, dass aus einer Menge mit mindestens 6 Elementen, 3 Elemente existieren, welche entweder alle miteinander verknüpft sind oder in gar keiner Relation stehen.


\section*{Aufgabe 4}
\subsection*{a)}

$w = CBBA$

$f_{\text{Alice}}(w)$ = 01101 0 0 01

$f_{\text{Bob}}(w)$ = 01 10 10 001

\subsection*{b)}
Voraussetzung für gleichheit der Decodierung ist unter anderem die gleichheit der Länge der decodierten Wörter.

Die Codierungen der Zeichen für Alice und Bob sind unterschiedlich lang: 
Länge $A_{\text{Alice}}$ = 2 Bit $A_{\text{Bob}}$ = 3 Bit

Länge $B_{\text{Alice}}$ = 1 Bit $B_{\text{Bob}}$ = 2 Bit

Länge $C_{\text{Alice}}$ = 5 Bit $C_{\text{Bob}}$ = 2 Bit

Länge $D_{\text{Alice}}$ = 2 Bit $D_{\text{Bob}}$ = 3 Bit

Daraus ergibt sich für ein Wort w mit der Länge n, wobei $n_{A}$, $n_{B}$, $n_{C}$, $n_{D}$ die Häufigkeit der jeweiligen Zeichen darstellen:

\begin{align*}
	\text{Länge}_{ALice} &= 2n_{A} + n_{B} + 5n_{C} + 2n_{D} \\
	\text{Länge}_{Bob} &= 3n_{A} + 2n_{B} + 2n_{C} + 3n_{D}
\end{align*}

Diese müssen gleich sein:


\begin{align*}
	2n_{A} + n_{B} + 5n_{C} + 2n_{D} &= 3n_{A} + 2n_{B} + 2n_{C} + 3n_{D} \\
	3n_{C}&= n_{A} + n_{B} + 3n_{D}
\end{align*}

\subsection*{c)}

Ein Wort ist "CEEBG".
\begin{align*}
	h_{Alice}(CEEBG) &= 01011011101101 \\
	h_{Bob}(CEEBG) &= 01011011101101
\end{align*}

A: Es kann kein Wort mit mit A geben für welches gilt: $h_{Alice}(w) = h_{Bob}(w)$, da $h_{Bob}(A)$ nicht durch $h_{Alice}(w)$ abgebildet werden kann.

B: Es gibt ein Wort: 

C: Es gibt ein Wort

D: Es kann kein Wort mit mit D geben für welches gilt: $h_{Alice}(w) = h_{Bob}(w)$, da $h_{Bob}(D)$ nicht durch $h_{Alice}(w)$ abgebildet werden kann.

E: Es gibt ein Wort.

F: Es kann kein Wort mit mit F geben für welches gilt: $h_{Alice}(w) = h_{Bob}(w)$, da $h_{Alice}(F)$ nicht durch $h_{Bob}(w)$ abgebildet werden kann.

G: Es gibt ein Wort

H: Es kann kein Wort mit mit H geben für welches gilt: $h_{Alice}(w) = h_{Bob}(w)$, da $h_{Bob}(H)$ nicht durch $h_{Alice}(w)$ abgebildet werden kann.

\end{document}

